\chapter{平面上的数量积}\label{ch:g11:scal}

数量积为一对向量对应一个\emph{数}，这个
数反映几何性质：它恰在两方向垂直时为零，
并可用于度量角度与长度。它是余弦定理
、圆的方程与垂直性计算背后的核心工具。
在空间中，同一工具是\cref{ch:g12:space}的研究对象。

\section{定义与初步计算}

\begin{definition}[数量积]\label{def:g11:scal:dot}
两个向量的\emph{数量积}\index{数量积}是
$\vec u$ 与 $\vec v$ 的一个数：
\[
\vec u \cdot \vec v
= \tfrac12\left(\norm{\vec u + \vec v}^2 - \norm{\vec u}^2
- \norm{\vec v}^2\right),
\]
其中 $\norm{\vec u}$ 为向量的\emph{模}\index{模}，即 $\vec u$ 的长度。
\end{definition}

\begin{theorem}[数量积的四种表示]\label{thm:g11:scal:faces}
设 $\vec u\,(x, y)$ 与 $\vec v\,(x', y')$ 为非零向量，且
$\theta$ 为它们之间的夹角。则
\begin{enumerate}
  \item \emph{（坐标）} $\vec u \cdot \vec v = xx' + yy'$；
  \item \emph{（余弦）} $\vec u \cdot \vec v =
        \norm{\vec u}\,\norm{\vec v}\cos\theta$；
  \item \emph{（投影）} $\vec u \cdot \vec v = \vec u \cdot \vec v'$，
        其中 $\vec v'$ 是 $\vec v$ 在
        $\vec u$ 方向上的正交投影；若 $\vec u = \vect{OA}$ 且
        $\vec v = \vect{OB}$，则
        $\vec u \cdot \vec v = \overline{OA} \times \overline{OH}$，即
        有向长度之积，其中 $H$ 是
        从 $B$ 向直线 $(OA)$ 所作垂线的垂足；
  \item \emph{（模）} 即\cref{def:g11:scal:dot}中的定义式。
\end{enumerate}
\end{theorem}

\begin{proof}[1 与 2 的证明]
\emph{1.} 在坐标下，$\norm{\vec u}^2 = x^2 + y^2$（毕达哥拉斯定理），
$\vec u + \vec v$ 的坐标为 $(x + x', y + y')$，于是
\begin{align*}
2\,\vec u \cdot \vec v
&= (x + x')^2 + (y + y')^2 - (x^2 + y^2) - (x'^2 + y'^2)\\
&= 2xx' + 2yy',
\end{align*}
展开平方并相消后即得。

\emph{2.} 选取坐标系，使 $\vec u$ 沿
正 $x$ 轴方向：$\vec u\,(\norm{\vec u}, 0)$，且
$\vec v\,(\norm{\vec v}\cos\theta,\ \norm{\vec v}\sin\theta)$
（见\cref{ch:g11:trigo}）。由公式 1 可得
$\vec u \cdot \vec v = \norm{\vec u}\,\norm{\vec v}\cos\theta + 0$。
公式 3 直接读出同一计算：$\vec v$ 在
$x$ 轴上的有向长度为 $\norm{\vec v}\cos\theta =
\overline{OH}$。
\end{proof}

\begin{omfigure}
\begin{tikzpicture}[scale=1.25]
  \coordinate (O) at (0,0);
  \coordinate (A) at (3.2,0);
  \coordinate (B) at (55:2.4);
  \coordinate (H) at (1.3767,0);
  \draw[-Stealth, omDef, thick] (O) -- (A)
    node[below, font=\small, pos=0.85] {$\vec u$};
  \draw[-Stealth, omThm, thick] (O) -- (B)
    node[above left, font=\small, pos=0.8] {$\vec v$};
  \draw[black, dashed, thin] (B) -- (H);
  \draw[thin] (H) ++(-0.18,0) -- ++(0,0.18) -- ++(0.18,0);
  \draw[omProp, very thick] (O) -- (H);
  \fill (H) circle (1.2pt) node[below, font=\small] {$H$};
  \node[below left, font=\small] at (O) {$O$};
  \draw[->,omThm,thin] (0.55,0) arc (0:55:0.55);
  \node[omThm, font=\small] at (27:0.78) {$\theta$};
\end{tikzpicture}

{\small 投影示意：$\vec u \cdot \vec v = \overline{OA} \times
\overline{OH}$，其中 $H$ 是从
$\vec v$ 的终点所作垂线的垂足。橙色线段 $OH$ 的有向长度为
$\norm{\vec v}\cos\theta$。}
\end{omfigure}

\begin{example}\label{ex:g11:scal:compute}
$\vec u\,(3, 1)$ 与 $\vec v\,(2, -4)$：
$\vec u \cdot \vec v = 6 - 4 = 2$。两向量的模分别为
$\norm{\vec u} = \sqrt{10}$，$\norm{\vec v} = \sqrt{20}$，于是
$\cos\theta = \frac{2}{\sqrt{10}\sqrt{20}} = \frac{2}{10\sqrt2}
= \frac{\sqrt2}{10}$：夹角略小于 $82^\circ$。
\end{example}

\begin{proposition}[代数运算规则]\label{prop:g11:scal:rules}
对所有向量及所有 $\lambda \in \R$：
\[
\vec u \cdot \vec v = \vec v \cdot \vec u, \qquad
\vec u \cdot (\vec v + \vec w) = \vec u \cdot \vec v + \vec u \cdot
\vec w, \qquad
(\lambda \vec u)\cdot \vec v = \lambda\,(\vec u \cdot \vec v),
\]
\[
\vec u \cdot \vec u = \norm{\vec u}^2,
\qquad
\norm{\vec u + \vec v}^2 = \norm{\vec u}^2 + 2\,\vec u\cdot\vec v +
\norm{\vec v}^2 .
\]
\end{proposition}

\begin{proof}
以上各式均可由坐标公式推出，例如
$\vec u \cdot (\vec v + \vec w) = x(x' + x'') + y(y' + y'') =
(xx' + yy') + (xx'' + yy'')$。最后一式展开
$\norm{\vec u + \vec v}^2 = (\vec u + \vec v)\cdot(\vec u + \vec v)$
利用前三条规则 --- 数量积的运算与普通乘积
完全相同，因此通常的代数恒等式对向量同样适用。
\end{proof}

\section{正交}

\begin{definition}[正交向量]\label{def:g11:scal:orthogonal}
若两个向量 \emph{正交}\index{正交向量}，记作
$\vec u \perp \vec v$，若 $\vec u \cdot \vec v = 0$。
\end{definition}

\begin{remark}
对非零向量，由\cref{thm:g11:scal:faces}的公式 2
可知 $\vec u \cdot \vec v = 0 \iff \cos\theta = 0 \iff$ 两向量的方向
垂直。零向量与任何向量都正交。
\end{remark}

\begin{example}
$\vec u\,(3, 2)$ 与 $\vec v\,(-4, 6)$：$\vec u \cdot \vec v = -12 + 12 =
0$：两向量正交。一般地，$\vec u\,(a, b)$ 总与
$\vec v\,(-b, a)$ 正交 --- 即一次四分之一周的旋转。
\end{example}

\section{余弦定理}

\begin{theorem}[余弦定理]\label{thm:g11:scal:cosines}
\index{余弦定理}
在三角形 $ABC$ 中，三边 $a = BC$、$b = CA$、$c = AB$，
角 $\widehat A$ 位于顶点 $A$ 处：
\[
a^2 = b^2 + c^2 - 2bc\,\cos\widehat A .
\]
\end{theorem}

\begin{proof}
记 $\vect{BC} = \vect{BA} + \vect{AC} = \vect{AC} - \vect{AB}$，并
利用\cref{prop:g11:scal:rules}展开：
\[
a^2 = \norm{\vect{BC}}^2
= \norm{\vect{AC}}^2 - 2\,\vect{AC}\cdot\vect{AB} + \norm{\vect{AB}}^2
= b^2 + c^2 - 2\,\vect{AB}\cdot\vect{AC}.
\]
由余弦公式，
$\vect{AB}\cdot\vect{AC} = c\,b\,\cos\widehat A$。
\end{proof}

\begin{remark}
当 $\widehat A$ 为直角时，$\cos\widehat A = 0$，余弦定理
退化为勾股定理 $a^2 = b^2 + c^2$：它是
毕达哥拉斯定理加上非直角情形的修正项。
\end{remark}

\begin{omfigure}
\begin{tikzpicture}[scale=1.15]
  \coordinate (A) at (0,0);
  \coordinate (B) at (4,0);
  \coordinate (C) at (1.2,2.2);
  \draw[omDef, thick] (A) -- (B) -- (C) -- cycle;
  \node[below left, font=\small] at (A) {$A$};
  \node[below right, font=\small] at (B) {$B$};
  \node[above, font=\small] at (C) {$C$};
  \node[below, font=\small] at (2,0) {$c$};
  \node[above left, font=\small] at (0.6,1.1) {$b$};
  \node[above right, font=\small] at (2.6,1.1) {$a$};
  \draw[->,omThm,thin] (0.7,0) arc (0:61.4:0.7);
  \node[omThm, font=\small] at (30:1.0) {$\widehat A$};
\end{tikzpicture}

{\small 余弦定理：已知两边 $b$、$c$ 及其夹角，
即可确定第三边 $a$。}
\end{omfigure}

\begin{example}\label{ex:g11:scal:cosines}
某三角形两边长分别为 $b = 5$、$c = 8$，它们的夹角
$\widehat A = \frac\pi3$。则
\[
a^2 = 25 + 64 - 2 \times 5 \times 8 \times \frac12 = 89 - 40 = 49,
\]
于是 $a = 7$。
\end{example}

\section{应用：圆与垂线}

\begin{proposition}[以已知线段为直径的圆]\label{prop:g11:scal:diameter}
点 $M$ 位于以 $[AB]$ 为直径的圆上，当且仅当
\[
\vect{MA} \cdot \vect{MB} = 0 .
\]
\end{proposition}

\begin{proof}
若 $M \neq A, B$，该条件说明 $M$ 处的角，
在三角形 $AMB$ 中为直角，而对线段张直角的点
恰为以该线段为直径的圆上的点。下面给出直接证明：
设 $O$ 为 $[AB]$ 的中点，$r = \frac{AB}{2}$。则
$\vect{MA} = \vect{MO} + \vect{OA}$，且
$\vect{MB} = \vect{MO} + \vect{OB} = \vect{MO} - \vect{OA}$，故
\[
\vect{MA}\cdot\vect{MB}
= \norm{\vect{MO}}^2 - \norm{\vect{OA}}^2 = MO^2 - r^2,
\]
该式恰在 $MO = r$ 时为零，即 $M$ 位于以 $O$ 为圆心、
$r$ 为半径的圆上。
\end{proof}

\begin{proposition}[直线的法向量]\label{prop:g11:scal:normal}
向量 $\vec n\,(a, b)$ 与直线
$d : ax + by + c = 0$ 的每一个方向向量都正交；称 $\vec n$ 为 $d$ 的 \emph{法向量}
\index{法向量}。两直线垂直，当且仅当
它们的法向量（或方向向量）正交。
\end{proposition}

\begin{proof}
$d$ 的一个方向向量为 $\vec u\,(-b, a)$
（见\cref{thm:g11:vect:cartesian}），且
$\vec n \cdot \vec u = a(-b) + b(a) = 0$。
\end{proof}

\begin{method}[使用法向量]\label{met:g11:scal:normal}
\begin{itemize}
  \item 过 $A$、法向量为 $\vec n\,(a,b)$ 的直线，可写为
        $\vect{AM} \cdot \vec n = 0$，展开后为 $ax + by + c = 0$，
        其中 $c$ 由 $A$ 确定。
  \item 两直线 $d: ax+by+c = 0$ 与 $d': a'x+b'y+c' = 0$ 垂直的判定：
        只需检验 $aa' + bb' = 0$。
\end{itemize}
\end{method}

\begin{example}\label{ex:g11:scal:perpline}
过 $P(2, -1)$ 且垂直于 $d : 3x - y + 4 = 0$ 的直线：$d$ 的一个方向向量为
$(1, 3)$，恰可作为所求直线的 \emph{法向量}。
设所求直线方程为 $x + 3y + c = 0$，将该点坐标代入得
$2 - 3 + c = 0$，故 $c = 1$，所求直线为 $x + 3y + 1 = 0$。
\end{example}

\begin{example}[圆的方程]\label{ex:g11:scal:circle}
圆心为 $\Omega(1, 2)$、半径为 $3$ 的圆，是所有
点 $M(x,y)$ 满足 $\Omega M^2 = 9$ 的集合，即：
\[
(x - 1)^2 + (y - 2)^2 = 9 .
\]
反之，$x^2 + y^2 - 2x - 4y - 4 = 0$ 可通过
配方（见\cref{ch:g11:quad}）化为 $(x-1)^2 + (y-2)^2 = 9$，
即为同一个圆。
\end{example}

\section{练习}

\begin{exercise}[$\star$]\label{exo:g11:scal:1}
计算下列各组向量的 $\vec u \cdot \vec v$：
\[
\vec u\,(2, 5),\ \vec v\,(3, -1); \qquad
\vec u\,(1, -4),\ \vec v\,(8, 2); \qquad
\norm{\vec u} = 3,\ \norm{\vec v} = 4,\ \theta = \frac{\pi}{3} .
\]
\end{exercise}

\begin{exercise}[$\star$]\label{exo:g11:scal:2}
$ABC$ 是边长为 $6$ 的正三角形。计算
$\vect{AB} \cdot \vect{AC}$ 与 $\vect{AB} \cdot \vect{BC}$。
\end{exercise}

\begin{exercise}[$\star$]\label{exo:g11:scal:3}
当 $t$ 取何值时，$\vec u\,(3, t)$ 与 $\vec v\,(t - 8, 1)$
互相正交？又当 $t$ 取何值时，$\vec u\,(t, 2)$ 与其减去
$\vec v\,(4, 2)$ 所得的向量正交，即 $\vec u \perp (\vec u - \vec v)$？
\end{exercise}

\begin{exercise}[$\star$]\label{exo:g11:scal:4}
计算 $\theta$，即 $\vec u\,(1, 2)$ 与 $\vec v\,(3, 1)$ 的夹角。
（先精确给出 $\cos\theta$，再将 $\theta$ 精确到整度）。
\end{exercise}

\begin{exercise}[$\star\star$]\label{exo:g11:scal:5}
一个三角形的边 $b = 4$、$c = 6$，两边的夹角为
$\widehat A = \frac{2\pi}{3}$，求第三边 $a$。
再在边长 $a = 7$、$b = 5$、$c = 4$ 的三角形中，求
$\cos\widehat A$，并判断 $\widehat A$ 是锐角还是钝角。
\end{exercise}

\begin{exercise}[$\star\star$]\label{exo:g11:scal:6}
设 $A(1, 1)$、$B(5, -1)$、$C(3, 5)$，计算
$\vect{AB} \cdot \vect{AC}$，$\widehat A$ 是否为直角？计算
三边长，并用余弦定理检验你的结论。
\end{exercise}

\begin{exercise}[$\star\star$]\label{exo:g11:scal:7}
写出圆心为 $\Omega(-2, 3)$、半径为 $5$ 的圆的方程。
再求下列圆的圆心与半径：
\[
x^2 + y^2 - 6x + 2y - 6 = 0 .
\]
\end{exercise}

\begin{exercise}[$\star\star$]\label{exo:g11:scal:8}
设 $A(-3, 0)$、$B(0, 2)$，利用\cref{prop:g11:scal:diameter}，求
以 $[AB]$ 为直径的圆的方程，并验证原点
$O(0,0)$ 在该圆上。
\end{exercise}

\begin{exercise}[$\star\star$]\label{exo:g11:scal:9}
求过 $P(1, 3)$ 且与
$d : 2x + y - 7 = 0$ 垂直的直线的方程，以及
垂足 $H$ 的坐标（即两直线的交点）。由此
求点 $P$ 到 $d$ 的距离。
\end{exercise}

\begin{exercise}[$\star\star$]\label{exo:g11:scal:10}
利用 $\norm{\vec u + \vec v}^2 = \norm{\vec u}^2 + 2\vec u\cdot\vec v +
\norm{\vec v}^2$ 及其对 $\vec u - \vec v$ 的类似等式，证明
\emph{平行四边形恒等式}：
\[
\norm{\vec u + \vec v}^2 + \norm{\vec u - \vec v}^2
= 2\norm{\vec u}^2 + 2\norm{\vec v}^2 ,
\]
并在平行四边形中解释该等式（对角线与边的关系）。
\end{exercise}

\begin{exercise}[$\star\star\star$]\label{exo:g11:scal:11}
设 $A$、$B$ 为两点，$AB = 4$，$I$ 为
$[AB]$ 的中点。
\begin{enumerate}
  \item 证明：对任意点 $M$，
        $\vect{MA}\cdot\vect{MB} = MI^2 - 4$。
        （插入点 $I$：$\vect{MA} = \vect{MI} + \vect{IA}$。）
  \item 由此求点 $M$ 的集合，使得
        $\vect{MA}\cdot\vect{MB} = 12$。
\end{enumerate}
\end{exercise}

\section{问题：数量积的两份礼物}

\begin{problem}[{周末问题 --- 两角和公式由一个数量积导出，
点到直线的距离由另一个数量积导出，
三条高线交于一点作为额外奖励}]
\label{pb:g11:scal:1}
数量积看上去只是记账式的运算 --- 相乘、相加 ---
却几乎不费吹灰之力地交出两件珍宝，它们
曾令数个世纪的角度推算束手无策：其一是三角学中的\emph{两角和公式}
（精确计算 $\cos 15^\circ$！），其二是
\emph{点到直线的距离}，可用一个简洁的分式表出。作为
临别礼物，它还证明了每个
三角形的三条高线交于一点。
\cref{thm:g11:scal:faces}的四个侧面都将登场。

\medskip
\noindent\textbf{第一部分 --- 熟练运算。}
\begin{enumerate}
  \item 对 $\vec u(3, 4)$ 与 $\vec v(-1, 2)$，计算
        $\vec u \cdot \vec v$、两向量的模，以及两向量
        的夹角（精确到十分之一度）。
  \item 求 $k$，使得 $(2, k)$ 与 $(3, -4)$ 正交
        （见\cref{def:g11:scal:orthogonal}）。
  \item 物理情境：一辆雪橇被拉动了 $10$ m，所用拉力为 $50$ N，
        与地面成 $60^\circ$ 角。计算所做的功
        $W = \vec F \cdot \vec d$。
  \item 余弦定理（见\cref{thm:g11:scal:cosines}）：两边长分别为
        $5$ 与 $7$，夹角为 $60^\circ$，
        精确求出第三边。
  \item 逆向：一个三角形的三边长为 $4$、$6$、$8$。求
        边 $8$ 所对角的余弦，判断该角的类型，
        并将该角精确到十分之一度。
\end{enumerate}

\medskip
\noindent\textbf{第二部分 --- 第一件礼物：两角和
公式。}
设 $\vec u = (\cos a, \sin a)$，$\vec v = (\cos b, \sin b)$：
两个单位向量，其方向角分别为 $a$ 与 $b$。
\begin{enumerate}[resume]
  \item 计算两次 $\vec u \cdot \vec v$ --- 一次用坐标，
        一次用模与两向量的夹角，
        并由此得出基本公式：
        \[
        \cos(a - b) = \cos a \cos b + \sin a \sin b .
        \]
  \item 以 $-b$ 代替 $b$（利用正弦与余弦的奇偶性！）
        求得 $\cos(a + b)$。
  \item 求 $\sin(a + b)$，需用余角关系式
        $\sin x = \cos\left(\frac\pi2 - x\right)$，
        取 $x = a + b$ 代入。
  \item 最后求精确值：计算 $\cos 15^\circ$（写作
        $\cos(45^\circ - 30^\circ)$）及 $\sin 75^\circ$，
        并用数值方法检验两者。
  \item 令 $b = a$，推出二倍角公式：
        $\cos 2a$（三种形式）及 $\sin 2a$。
  \item 由 $\cos 2a = 1 - 2\sin^2 a$，求
        $\sin^2 15^\circ$，再将 $\sin 15^\circ$ 表示为
        复合根式 --- 并验证其平方与
        第 9 题的方法所得的值 $\frac{\sqrt6 - \sqrt2}{4}$
        相符。
\end{enumerate}

\medskip
\noindent\textbf{第三部分 --- 第二件礼物：点到直线的
距离。}
\begin{enumerate}[resume]
  \item 设直线 $d$ 为 $ax + by + c = 0$，其法向量为
        $\vec n(a, b)$（见\cref{prop:g11:scal:normal}），
        又设 $P(x_0, y_0)$ 为一点，$A$ 为 $d$ 上任一点，
        说明点 $P$ 到 $d$ 的距离为
        $\dfrac{\abs{\vect{AP} \cdot \vec n}}{\norm{\vec n}}$,
        并由此推出公式
        \[
        \operatorname{dist}(P, d)
        = \frac{\abs{a x_0 + b y_0 + c}}{\sqrt{a^2 + b^2}} .
        \]
        求点 $P(5, 6)$ 到直线的距离，直线为
        $3x + 4y - 12 = 0$.
  \item 写出以 $(5, 6)$ 为圆心、
        且与该直线相切的圆的方程。
  \item 检验公式：计算原点到
        直线 $3x + 4y - 12 = 0$ 的距离，然后通过
        求出从原点所作垂线的垂足
        来加以验证。
  \item 三角形 $A(0,0)$、$B(6,0)$、$C(2,5)$：
        以底边 $[AB]$ 计算其面积；再求点
        $B$ 到直线 $(AC)$ 的距离，并验证由此得出的面积
        与前者相同。
  \item 圆心为 $(2, 1)$、半径为 $2$ 的圆，与
        直线 $x + y - 6 = 0$：计算圆心
        到该直线的距离，并判断二者的位置关系（割线、
        切线、相离）。
\end{enumerate}

\medskip
\noindent\textbf{第四部分 --- 额外奖励：三条高线共点。}
\begin{enumerate}[resume]
  \item 推广\cref{exo:g11:scal:11}：对任意两点
        $A$、$B$，设其中点为 $I$，证明
        \[
        \vect{MA} \cdot \vect{MB} = MI^2 - \frac{AB^2}{4}
        \quad\text{对每一点 } M .
        \]
  \item 用一行推出\cref{prop:g11:scal:diameter}：
        点 $M$ 的集合，满足
        $\vect{MA} \cdot \vect{MB} = 0$，即以
        $[AB]$ 为直径的圆。
  \item 三条高线：设 $H$ 为过 $A$ 与过 $B$ 的两条高线的交点，
        其中三角形为 $ABC$，
        （故 $\vect{HA} \cdot \vect{BC} = 0$ 且
        $\vect{HB} \cdot \vect{CA} = 0$）。证明恒等式
        \[
        \vect{HA} \cdot \vect{BC} + \vect{HB} \cdot \vect{CA}
        + \vect{HC} \cdot \vect{AB} = 0
        \]
        --- 从 $H$ 出发、借助沙勒关系把一切向量展开 --- 并
        由此得出 $H$ 也在过 $C$ 的高线上，
        即三条高线共点（交点为
        \emph{垂心}\index{垂心}）。
  \item 收尾：四个侧面，三座奖杯。在\cref{thm:g11:scal:faces}中，
        哪一个侧面承载了两角和公式，
        哪一个承载距离公式，哪一个承载垂心 ---
        为何“两次计算同一个量”是贯穿全章的
        座右铭？
\end{enumerate}
\end{problem}
